% Generated draft: requires mathematical and source review.
\documentclass[11pt]{article}
\usepackage[margin=25mm]{geometry}
\usepackage{fontspec}
\setmainfont{Latin Modern Roman}
\usepackage{amsmath,amssymb,mathtools,unicode-math}
\setmathfont{Latin Modern Math}
\usepackage{enumitem,needspace,xurl,hyperref}
\setlength{\parindent}{0pt}
\setlength{\parskip}{5pt}
\setlength{\emergencystretch}{3em}
\newcommand{\N}{\mathbb{N}}
\newcommand{\Z}{\mathbb{Z}}
\newcommand{\Q}{\mathbb{Q}}
\newcommand{\R}{\mathbb{R}}
\newcommand{\C}{\mathbb{C}}
\newcommand{\F}{\mathbb{F}}
\newcommand{\A}{\mathbb{A}}
\newcommand{\PP}{\mathbb{P}}
\newcommand{\E}{\mathbb{E}}
\newcommand{\eps}{\varepsilon}
\newcommand{\dd}{\,\mathrm{d}}
\begin{document}
\section*{UnsolvedMath ID: 10400103}
\textbf{Source problem code:} catalogue entry
\subsection{Definitions and mathematical statement}
A quandle is a set \(X\) with an idempotent operation \(*\), bijective right translations, and right self-distributivity \((x*y)*z=(x*z)*(y*z)\). In a usual knot-cocycle construction, take finite \(X\), an abelian coefficient group \(A\) written multiplicatively, and a normalized quandle \(2\)-cocycle \(\theta:X^2\to A\), satisfying
\[
 \theta(x,x)=1,\qquad
 \theta(x,y)\theta(x*y,z)=\theta(x,z)\theta(x*z,y*z).
\]
For an oriented knot diagram \(D\), color its arcs by \(X\) subject to the quandle crossing relation. With the compatible choice of under-arc and over-arc colors \(x_c,y_c\) and crossing sign \(\eps_c\), the usual cocycle state sum has form
\[
 \Phi_{X,\theta}(K)=
 \sum_{\mathcal C\text{ an }X\text{-coloring of }D}
      \prod_{c\text{ a crossing}}\theta(x_c,y_c)^{\eps_c}
       \ \in\Z[A].
\]
The negative-crossing color convention and any normalization must be fixed to obtain the intended invariant. This supplies a common example of the named family, not a source-selected choice of \(X,A,\theta\).

Problem 5.7 asks to find relations between quandle cocycle invariants and previously known knot invariants, explicitly including quantum invariants. The record does not select particular quandles or cocycles, a quantum group or representation/category, or a relation to prove. Equalities, specializations, reconstruction formulas, and comparisons of distinguishing power would be different concrete targets, so none is silently imposed. The general comparison program, including quantum invariants as examples rather than the sole permitted family, is the supplied mathematical request.
\subsection{Short English statement}
Relate quandle-cocycle knot invariants to other known knot invariants, including quantum invariants. The source gives no specific identity or target family. Its supplied partial-progress assessment says some relations are known, without listing them or giving references.
\subsection{Sources}
\begin{itemize}
\item Local UnsolvedMath registry, permanent ID \(10400103\), exact source code \texttt{catalogue entry}, Problem 5.7. No individual proposer or proposed year is supplied.
\item Ohtsuki, \emph{Problems on invariants of knots and 3-manifolds} (2002), Problem 5.7, PDF page 91: \url{https://msp.org/gtm/2002/04/gtm-2002-04-024s.pdf}. Bold headings distinguish 231 numbered items from surrounding references.
\item Extraction pdf-font-structure, accessed 2026-07-29. Source evidence NEEDS\_REVIEW presents an open problem and requires a current primary-literature audit. Registry status partially\_solved and research classification PARTIAL-PROGRESS report some known relations between cocycle and quantum invariants, without specific primary citations.
\item Dataset attribution: Ulam AI / UnsolvedMath Contributors, CC BY 4.0. Source rights NEEDS\_REVIEW: publicly accessible author/publisher copy, redistribution terms requiring review.
\end{itemize}
\end{document}
